\documentclass{article}
\input{../../../lib.sty}

\title{\texttt{fn new\_negative\_binomial\_rdp\_curve}}
\author{Michael Shoemate}

\begin{document}
\maketitle

Proves soundness of \texttt{fn new\_negative\_binomial\_rdp\_curve} in \texttt{mod.rs}.

\section{Hoare Triple}
\subsection*{Precondition}
\texttt{base\_curve} upper-bounds the base Rényi divergence at every order $> 1$,
\texttt{eta} $\ge 0$, \texttt{mean} $> 0$ is finite, and
\texttt{log\_inv\_gamma\_hi\_from\_x(x)} succeeds whenever it is called.

\subsection*{Pseudocode}
\lstinputlisting[language=Python,firstline=2,escapechar=|]{./pseudocode/new_negative_binomial_rdp_curve.py}

\subsection*{Postcondition}
\begin{theorem}
\label{postcondition}
For every $\alpha > 1$, the returned function either returns \texttt{Err(e)} due to a propagated
helper error, or returns an upper bound on the implemented Theorem 2 expression
\[
\epsilon(\alpha)
\;+\;
(1+\eta)\!\left(1 - \frac{1}{\alpha}\right)\epsilon(\alpha)
\;+\;
\frac{(1+\eta)\log(1/\gamma)}{\alpha}
\;+\;
\frac{\log(\mathbb{E}[K])}{\alpha - 1},
\]
where $\epsilon(\cdot)$ is the base Rényi divergence curve, $\gamma = 1 - e^{-x}$, and the input
\texttt{mean} upper-bounds $\mathbb{E}[K]$.
\end{theorem}

\section{Proof}
Fix $\alpha > 1$.

\begin{lemma}
\label{guard}
If $\alpha \le 1$, the returned value is $+\infty$, which is a valid upper bound.
\end{lemma}

\begin{proof}
The expression is finite or infinite, and $+\infty$ upper-bounds either case.
\end{proof}

\begin{lemma}
\label{main}
If $\alpha > 1$ and the function returns \texttt{Ok(out)}, then \texttt{out} upper-bounds the
displayed expression.
\end{lemma}

\begin{proof}
Let $\epsilon(\cdot)$ denote the exact base Rényi divergence curve. By precondition,
\[
\texttt{eps\_alpha} \ge \epsilon(\alpha).
\]
Also \texttt{one\_plus\_eta\_hi} $\ge 1 + \eta$, \texttt{inv\_alpha\_lo} $\le 1/\alpha$, and hence
\[
\texttt{one\_minus\_inv\_alpha\_hi} \ge 1 - \frac{1}{\alpha}.
\]
All of these quantities are nonnegative, so upward-rounded multiplication yields
\[
\texttt{term1\_hi} \ge (1+\eta)\!\left(1 - \frac{1}{\alpha}\right)\epsilon(\alpha).
\]
By the proof of \texttt{log\_inv\_gamma\_hi\_from\_x},
\[
\texttt{log\_inv\_gamma\_hi} \ge \log(1/\gamma).
\]
Since \texttt{alpha} is a positive exact denominator,
\[
\texttt{term2\_hi} \ge \frac{(1+\eta)\log(1/\gamma)}{\alpha}.
\]
By the proof of \texttt{log\_mean\_hi},
\[
\texttt{log\_mean\_hi\_} \ge \log(\texttt{mean}) \ge \log(\mathbb{E}[K]).
\]
Also \texttt{alpha\_minus\_one\_lo} $\le \alpha - 1$ with a positive denominator, so
\[
\texttt{term3\_hi} \ge \frac{\log(\mathbb{E}[K])}{\alpha - 1}.
\]
Finally, \texttt{out} is formed by upward-rounded addition of
\texttt{eps\_alpha}, \texttt{term1\_hi}, \texttt{term2\_hi}, and \texttt{term3\_hi}, so it
upper-bounds the full displayed expression.
\end{proof}

\begin{proof}[Proof of Theorem~\ref{postcondition}]
By Lemma~\ref{guard} for $\alpha \le 1$, and Lemma~\ref{main} otherwise.
\end{proof}

\end{document}
